% Research contribution D: files and identifiers owned by cyclotomic_coverings.
\subsection{Cyclotomic points: abelianization and explicit finite covers}

\paragraph{Statement D1 (Proved).}\label{res:D1}
Let $F$ be a finite extension of $\mathbb Q(\mu_\infty)$, let
$U=X\setminus S$ be a smooth geometrically connected hyperbolic curve over
$F$, and let $x,y\in U(F)$.  We use geometric conjugacy classes of continuous
sections and the natural continuous quotient
$E_U\twoheadrightarrow E_U/\overline{[\pi,\pi]}$.
For $g(X)\geq1$, equality of the images of the two point sections under this
quotient implies $x=y$.  For $g(X)=0$ and $x\ne y$, equality occurs precisely
when every geometric point of $S$ is $F$-rational and
\[
 \frac{h(a)}{h(b)}\in\mu_\infty\qquad(a,b\in S),\qquad
 h(t)=\frac{t-x}{t-y}.
\]
Here we choose any $F$-coordinate on $X\simeq\mathbb P^1$; the displayed
ratios have the evident homogeneous interpretation at infinity and do not
depend on the coordinate.  In particular, on
$U=\mathbb P^1\setminus\{0,1,\infty\}$ equality is equivalent to
$x/y,(1-x)/(1-y)\in\mu_\infty$.  Its distinct ordered pairs are exactly
\[
 (x,y)=\left(\frac{a(1-b)}{a-b},\frac{1-b}{a-b}\right),\quad
 a,b\in\mu_\infty,\quad a,b\ne1,\quad a\ne b.
\]
This answers the abelianized question of scan 62 (manuscript p.~30).

\begin{proof}
We make the comparison with Kummer theory explicit.  Choosing $y$ as origin
of the generalized Abel--Jacobi map gives a morphism $U\to J_S$.  Its map on
geometric abelian fundamental groups is the identification with the full
Tate module $T(J_S)$: at level $n$, pulling back multiplication by $n$ on
$J_S$ gives the abelian covering with monodromy
$\pi\twoheadrightarrow\pi^{\rm ab}/n$.  This assertion follows alternatively
from the degree-zero divisor description of $J_S$: the compact part is
$T(J_X)$ and each puncture contributes its tame inertia generator, with the
single sum relation.  Both constructions give the exact sequence
\[
 0\longrightarrow
 \bigl(\widehat{\mathbb Z}(1)[S]/\widehat{\mathbb Z}(1)\bigr)
 \longrightarrow \pi^{\rm ab}\longrightarrow T(J_X)\longrightarrow0.
\]
The arithmetic comparison sends the difference of point sections to the
Kummer class of $[x-y]\in J_S(F)$.  A Kummer class vanishes precisely when
this element belongs to every $nJ_S(F)$: the finite-level statement is the
exact sequence for $[n]:J_S\to J_S$, and compatible choices follow from the
inverse system of finite nonempty division-point sets.  Thus the comparison
is proved in both directions.

For positive genus its image in $J_X(F)$ is divisible.  The arithmetic
input of Part A says that this image is zero.  A principal divisor $x-y$
with $x\ne y$ would give a degree-one morphism $X\to\mathbb P^1$; such a
morphism between smooth proper curves is an isomorphism, contrary to
$g(X)\geq1$.

In genus zero the existence of $x$ identifies $X$ with $\mathbb P^1$.
The divisor description gives
$J_S=\operatorname{Res}_{S/F}\mathbb G_m/\mathbb G_m$ and
$J_S(F)=\mathcal O(S)^*/F^*$; the equality on rational points follows also
from $H^1(F,\mathbb G_m)=0$.  The class of $x-y$ is represented by the tuple
$(h(a))_{a\in S}$, up to inversion, which is irrelevant for divisibility.
Take a finite Galois splitting field $L/F$ of $S$.  After restriction to $L$,
divisibility and Part A(iv) show that $h(a)/h(b)$ is a root of unity for
all $a,b$.  These roots of unity already belong to $F$.  Write
$h(a)=t\zeta_a$ with $\zeta_a\in\mu_\infty$.  Equivariance implies that
$\lambda(\sigma)=\sigma(t)/t$ is a character
$G=\operatorname{Gal}(L/F)\to\mu_\infty$.
Let $n=|G|$.  Divisibility over $F$ supplies an equivariant tuple $(b_a)$
and $c\in F^*$ with $b_a^n=c h(a)$.  Then $b_a/b_b$ is a root of unity,
so the same argument supplies a character $\psi:G\to\mu_\infty$ with
$\psi^n=\lambda$.  Every character of $G$ has exponent dividing $n$.
Consequently $\lambda=1$, all $h(a)$ belong to $F$, and the inverse
$F$-M\"obius transformation shows $a\in F$ for every cusp.  Conversely,
when $S$ is split and these ratios are roots of unity, every coordinate of
the class in the split torus has roots of every order in $F$, so the class
is divisible.  The parametrization for three punctures follows by solving
$x=ay$, $1-x=b(1-y)$; the listed exclusions are exactly those ensuring
$x,y\notin\{0,1\}$ and $x\ne y$.
\end{proof}

% Further statements are appended below the cyclotomic-base discussion.

\paragraph{Statement D2 (Proved).}\label{res:D2}
For every finite extension $F'/F$ and every hyperbolic curve $U/F'$, the
map from $U(F')$ to geometric conjugacy classes of full continuous sections
is injective.  The same holds over every finitely generated extension of
$F$.  These assertions make no assertion of surjectivity.

\begin{proof}
We give the finite-cover comparison, including its descent.  Suppose two
point sections have the same class.  Conjugate the chosen lift of the
second geometric point so that their actual homomorphisms
$s:G_{F'}\to E_U$ agree.  Choose a finite-index characteristic subgroup
$H\subseteq\pi$ whose covering has compactification of positive genus.
Such a subgroup exists: for positive genus take $H=\pi$; in genus zero
choose three punctures over the algebraic closure and pull back the
Fermat cubic cover $w^3+v^3=1$ by a M\"obius coordinate carrying them to
$0,1,\infty$.  Removing the other punctures leaves a connected finite
\'etale covering of $U_{\overline F'}$ of genus one.  Taking its Galois
closure and then the intersection of all subgroups of the resulting
index produces an open characteristic subgroup contained in the original
one.  There are finitely many such subgroups because $\pi$ is finitely
generated.  The induced finite map of compactifications has positive-genus
target, and Riemann--Hurwitz gives positive genus for its source.

The open subgroup $H s(G_{F'})\subseteq E_U$ defines a connected finite
\'etale covering $V\to U$ over $F'$.  It has geometrically connected fibre
because its projection to $G_{F'}$ is surjective and its geometric subgroup
is exactly $H$.  Both geometric lifts chosen above are fixed by the common
$s(G_{F'})$, so descend to $F'$-points $\widetilde x,\widetilde y\in V(F')$.
Their sections in $E_V=H s(G_{F'})$ are the same actual homomorphism.
The morphism from $V$ to its smooth proper compactification carries them
to points with the same abelianized section.  Statement D1 in positive
genus gives $\widetilde x=\widetilde y$, hence $x=y$.

For finitely generated extensions we use the same argument and the AVKF
property established in Part A/E.  The exact external statements are
Tsujimura \cite[Proposition 1.2(ii), Proposition 4.2]{D_Tsujimura2022}:
finite generation preserves AVKF, and geometric points of a hyperbolic
curve over an AVKF field are separated by their germs of point sections.
Hoshi--Mochizuki--Tsujimura \cite[Definition 3.3(vi), Proposition 6.4]{D_HMT}
records the same injectivity for AVKF subfields of $\overline{\mathbb Q}$.
We have supplied the comparison rather than assuming section surjectivity.
\end{proof}

\paragraph{Statement D3 (Proved).}\label{res:D3}
Put $F=\mathbb Q(\mu_\infty)$ and $\xi=\zeta_6$.  For all positive
integers $n,m$ and every $\zeta\in\mu_n$, the finite \'etale covering of
$U_{0,3}$ defined by
\[
 w^n=z,\qquad u^m=1-\zeta w
\]
has isomorphic $G_F$-sets as fibres over $\xi$ and $\xi^{-1}$.
The bound $1\leq n,m\leq12$, with every $\zeta\in\mu_n$, is checked
exactly in \texttt{code/D\_coverings.py}; the argument below proves the
unbounded assertion.  The restriction $\zeta^n=1$ ensures that its extra
branch value is $z=1$, so that this is indeed a covering of the specified
three-punctured line.

\begin{proof}
Every $w$ in either fibre is a root of unity, hence belongs to $F$.
The map $w\mapsto w'=\zeta^{-2}w^{-1}$ bijects the first fibre of $w^n=z$
with the second.  We have the exact identity
\[
 1-\zeta w'=-\frac{1-\zeta w}{\zeta w}.
\]
Choose $\lambda_w\in\mu_\infty\subset F$ with
$\lambda_w^m=-(\zeta w)^{-1}$.  Then
$(w,u)\mapsto(w',\lambda_w u)$ is a $G_F$-equivariant bijection of the
finite fibres.  It commutes with the second-stage $\mu_m$-action.  This
comparison applies to these covers; the cover in D4 supplies a further
finite Galois-set invariant.
\end{proof}

\paragraph{Statement D4 (Proved).}\label{res:D4}
The full point sections at $\zeta_6$ and $\zeta_6^{-1}$ over
$F=\mathbb Q(\mu_\infty)$ are separated by an explicit connected finite
\'etale covering of degree $81$.  In particular, the smallest instance
specified in Part D(c) is decided.

\begin{proof}
Let $E$ be the smooth Fermat cubic $w^3+v^3=1$, with the flex
$(1:-1:0)$ as origin, and let $\varphi:E\to\mathbb P^1$ be $z=w^3$.
It has degree $9$ and ramifies only above $0,1,\infty$.  The formulas
\[
 X=\frac{12}{w+v},\qquad Y=\frac{36(w-v)}{w+v}
\]
identify its smooth projective model with $Y^2=X^3-432$; their inverse is
$w=(36+Y)/(6X)$, $v=(36-Y)/(6X)$.  Fix $a=\zeta_{18}$,
$\omega=a^6$, and $P=(a,a^{-1})\in E(F)$ in Fermat coordinates.
Our program performs the group law in the exact field
$\mathbb Q[a]/(a^6-a^3+1)$ and verifies
\[
 3P=(0,24a^3-12)\ne O,\qquad9P=O.
\]
The nine points above $\xi^{-1}$ are
$Q_{ij}=(\omega^i a^{-1},\omega^j a)$, $0\leq i,j<3$.
For each of them it verifies
\[
 3(Q_{ij}-P)=3P,\qquad9(Q_{ij}-P)=O.
\]
Thus all nine differences have exact order $9$.  Above $\xi$, the
corresponding differences have orders $1,3,\ldots,3$.
These are finite identities checked using rational arithmetic; the output
prints every pair $(i,j)$ and every triple.

We next prove the arithmetic bound that turns these identities into a
Galois-set separation.  The automorphism
$(X,Y)\mapsto(\omega X,Y)$ gives an action of
$\mathbb Z[\omega]$ on $E$, defined over $\mathbb Q(\omega)$.
Its action on $T_3E$ is faithful
\cite[II, Proposition 6.7]{D_MilneEC}.
The $3$-adic Tate module is a free rank-one module over
$\mathbb Z_3[\omega]$: it is free of rank two over $\mathbb Z_3$,
$T^2+T+1$ is irreducible over $\mathbb Q_3$, and
$\mathbb Z_3[\omega]$ is the valuation ring of the ramified quadratic
extension $\mathbb Q_3(\omega)$ (the element $\omega-1$ satisfies the
Eisenstein polynomial $T^2+3T+3$).  The module is finitely generated and
torsion-free over this valuation ring, hence free; comparison of
$\mathbb Z_3$-ranks gives rank one.  For every $n\geq1$ the natural map
\[
 T_3E/3^nT_3E\ \xrightarrow{\sim}\ E[3^n](\overline{\mathbb Q})
\]
is equivariant both for Galois and for the endomorphisms just specified.

The displayed Weierstrass equation has discriminant $-2^{12}3^9$, so
the elliptic curve has good reduction at $7$, and the program counts
$\#E(\mathbb F_7)=9$.  We use the following precise source statements:
Milne \cite[II, Proposition 6.3; V, Propositions 8.1 and 8.3]{D_MilneEC}
gives the rank-two prime-to-characteristic torsion module, unramifiedness
of the Tate representation at good primes away from its prime, and
$\operatorname{tr}(\operatorname{Frob}_p)=p+1-\#E(\mathbb F_p)$;
II, Proposition 6.4 identifies the Tate determinant with the degree of
an elliptic endomorphism.  Hence a Frobenius lift at $7$ acts on $T_3E$
with characteristic polynomial $T^2+T+7$.  As $7$ splits in
$\mathbb Q(\omega)$, a Frobenius lift $g$ in a decomposition group at $7$
restricts to the identity on $\mathbb Q(\omega)$.  It therefore commutes
with the specified endomorphism $\omega$ and, on the rank-one
$\mathbb Z_3[\omega]$-module, is multiplication by
$\alpha=1+3\omega$ or its conjugate.
Complex conjugation $c$ conjugates $\omega$ to $\bar\omega$, so the
commutator $cgc^{-1}g^{-1}$ acts by
$\bar\alpha/\alpha$ (or its inverse).  Each cyclotomic extension of
$\mathbb Q$ is abelian, so this commutator fixes $F$ pointwise.  It
therefore fixes every point defined over $F$.  Every $3$-power torsion point
$R\in E(F)$ belongs to some $E[3^n]$; the equivariant quotient above
therefore gives $(\bar\alpha/\alpha-1)R=0$, with the same conclusion
for the inverse scalar.  Since $\alpha$ is a unit in
$\mathbb Z_3[\omega]$ (its norm is $7$), this implies
$(\alpha-\bar\alpha)R=0$.
Since
\[
 (\alpha-\bar\alpha)(-1-2\omega)
 =(3+6\omega)(-1-2\omega)=9,
\]
we obtain $9R=0$.  The program checks the multiplication matrix
$\left(\begin{smallmatrix}3&-6\\6&-3\end{smallmatrix}\right)$,
its Smith form $(3,9)$ and the kernels modulo $9,27,81$.

Now let $M:E\to E$ be $R\mapsto[3]R+P$ and set
$W=M^{-1}(\varphi^{-1}(U_{0,3}))$.  The composite
$W\xrightarrow{M}\varphi^{-1}(U_{0,3})\xrightarrow{\varphi}U_{0,3}$
is defined over $F$: the curve and $\varphi$ are defined over $\mathbb Q$,
and $P\in E(\mathbb Q(\zeta_{18}))\subset E(F)$ defines the translation.
Over $U_{0,3}$, the map $\varphi$ is given by
$w^3=z$, $v^3=1-z$, with $w,v$ invertible; the two monic Kummer
equations give a finite \'etale algebra of rank $9$.
Multiplication by $3$ is finite \'etale of degree $9$ in characteristic
zero, and translation is an isomorphism, so the displayed composite is
finite \'etale of degree $81$.  Moreover $W$ is a nonempty open
subscheme of the geometrically integral elliptic curve $E$, hence is
geometrically connected.  Its fibre over $\xi$ contains the $F$-point
$R=O$, because $M(O)=P$ and $\varphi(P)=a^3=\xi$.
An $F$-point over $\xi^{-1}$ would satisfy
$3R=Q_{ij}-P$ for some $(i,j)$, forcing $R$ to have exact order $27$,
since $9R=3P\ne O$ and $27R=O$.  This contradicts the proved bound.
The two finite $G_F$-sets therefore have
different numbers of fixed elements.  Conjugate full sections induce
isomorphic $G_F$-sets on every finite \'etale fibre; this one disproves
conjugacy.  This last implication is the fibre-functor map from the
arithmetic fundamental group to the permutation group of the finite
covering, restricted to each actual section.
\end{proof}

\subsection{The proposed extension of the manuscript to cyclotomic bases}

\paragraph{Statement D5 (Refuted, with explicit counterexample).}\label{res:D5}
The extension of Conjecture A to $F=\mathbb Q(\mu_\infty)$ is false,
even for proper hyperbolic curves.  The smooth projective normalization
$C/F$ of
\[
                         t^6+2y^3+4=0
\]
has genus $4$, has no $F$-rational point, and admits a continuous section
of $E_C\twoheadrightarrow G_F$.  The same counterexample persists after
adjoining any finite number of independent transcendental variables to
$F$.  Thus ``tous nos beaux \'enonc\'es'' (scan 62, manuscript p.~30) cannot include
the proposed surjectivity statement.  The point-injectivity and centralizer
statements proved separately remain valid.

\begin{proof}
The morphism $(t,y)\mapsto(t^2:y:1)$ maps the affine equation to
$D:X^3+2Y^3+4Z^3=0$ and extends to a morphism of smooth projective curves.
For $L=\mathbb Q(\mu_N)$ and a place above $2$, the ramification index
$e$ is a power of $2$: the odd-order roots of unity give unramified local
extensions, while the $2$-power roots give ramification degree
$\varphi(2^r)$.  In the integer-normalized valuation of $L$, the three
nonzero summands in the cubic have respective residues $0,e,2e$ modulo
$3$.  These residues are distinct.  Hence the term of least valuation is
unique and their sum cannot vanish.  Zero coordinates do not change the
argument for the remaining terms.  Thus $D(L)=\varnothing$ for every
$N$, and $C(F)=\varnothing$.  This is the case $p=2$, $a=b=c=1$ of
Clark \cite[Theorem 1]{D_Clark2006}, whose full hypothesis is
$p\equiv-1\pmod3$ and $p\nmid abc$; we have reproduced the valuation
proof for the example used here.

The function-field extension $F(t)\subset F(t,y)$ has degree three:
$-(t^6+4)/2$ has a simple zero and so is not a cube even over
$\overline F(t)$.  The six roots of $t^6+4$ are simple.  They are the six
branch values of index three; the pole at infinity has order six, so
infinity has three unramified preimages.  Riemann--Hurwitz gives
$2g-2=3(-2)+6(3-1)=6$, hence $g=4$.  The exact polynomial and valuation
residue checks are in \texttt{code/D\_cyclotomic\_counterexample.py}.

We give the section-existence argument, including the finite lifting
step, to avoid assuming any section conjecture.  First,
$\operatorname{Br}(L)=0$ for every finite extension $L/F$.
A central simple algebra over $L$ descends to a number field $k\subset L$.
Milne \cite[VII, Proposition 7.2]{D_MilneCFT} states that every Brauer
class of a number field is killed by a finite cyclic cyclotomic extension.
That extension of $k$ is contained in $L$, which contains all roots of
unity.  For clarity, the local-global input behind that proposition is
VIII, Theorem 4.2: the map from the Brauer group of a number field to the
direct sum of its local Brauer groups is injective.  III, Theorem 2.1
identifies the nonarchimedean local group with $\mathbb Q/\mathbb Z$ and
restriction with multiplication by the local degree.  Cyclotomic
extensions provide arbitrarily divisible unramified local degrees,
through $\mu_{p^d-1}$ at $p$; a finite compositum kills the finitely many
nonzero local invariants, and adjoining $i$ kills real invariants.

Next $H^2(G_F,A)=0$ for every finite continuous $G_F$-module $A$.
Here is the group-theoretic reduction.  It suffices to treat its
$\ell$-primary parts.  Let $P$ be a pro-$\ell$ Sylow subgroup of $G_F$.
For every open $V$ containing $P$, Kummer theory and the preceding
Brauer vanishing give $H^2(V,\mathbb F_\ell)=0$, since $\mu_\ell\subset F$.
Continuous cochains give
$H^2(P,\mathbb F_\ell)=\varinjlim_{P\subseteq V}H^2(V,\mathbb F_\ell)$:
a locally constant cochain factors through a finite quotient, and its
cocycle or coboundary identities extend to an open neighbourhood $PN$ of
$P$ for a sufficiently small open normal $N$.  A finite $P$-module has a
filtration with trivial $\mathbb F_\ell$ factors: reduce by powers of
$\ell$, and use the nonzero invariants of every finite $\ell$-group on a
nonzero vector space in characteristic $\ell$.  The latter follows by
induction from a central element of order $\ell$, whose operator $g-1$
is nilpotent.  The cohomology long exact sequences now give
$H^2(P,A)=0$.  Restriction to $P$ is injective on the $\ell$-primary
cohomology: vanishing on $P$ gives vanishing on some such $V$, and
corestriction composed with restriction multiplies by $[G_F:V]$, prime
to $\ell$.  The transfer identity is Milne
\cite[II, Proposition 1.30]{D_MilneCFT}; its profinite version follows
from the same finite-quotient cochains.  Thus the claimed vanishing holds.

We now solve a finite embedding problem $G_F\to A_0\leftarrow B_0$
with $B_0\twoheadrightarrow A_0$ surjective.  Replace $A_0$ by the
image of the first map and $B_0$ by its preimage.  Choose a subgroup
$H\subseteq B_0$ minimal among those surjecting onto $A_0$.  Its kernel
$N$ lies in the Frattini subgroup of $H$: a maximal subgroup not
containing $N$ would itself surject onto $A_0$.  Moreover $N$ is
nilpotent.  Indeed, for any Sylow subgroup $P_N$ of the normal subgroup
$N$, the finite Frattini argument gives
$H=N_H(P_N)N$; the Frattini property forces $N_H(P_N)=H$.
All Sylow subgroups of $N$ are therefore normal.  When $N\ne1$, its
centre contains a nonzero elementary abelian subgroup invariant under
$H$, and hence a minimal nonzero $H$-normal elementary abelian subgroup
$D$.  Induction on $|N|$ solves the problem for $H/D\to A_0$.
Pulling back $H\to H/D$ by that solution gives an extension of $G_F$
by the finite abelian module $D$.  Its class in $H^2(G_F,D)$ vanishes,
so it has a continuous splitting, yielding the required weak solution.
The last statement follows directly from a continuous set-theoretic
section and its two-cocycle; such a section exists for a finite-kernel
map of profinite groups by using finitely many clopen cosets.

Every profinite extension of $G_F$ with finite kernel consequently
splits.  In detail, choose an open normal subgroup $V$ of the extension
group meeting its finite kernel trivially.  The extension is the fibre
product of the finite map $E/V\to G_F/p(V)$ with
$G_F\to G_F/p(V)$, so the solved finite embedding problem supplies its
splitting.  Finally the geometric fundamental group of $C$ is finitely
generated and has a countable cofinal sequence
$\pi=N_0\supset N_1\supset\cdots$ of open characteristic subgroups.
Starting with the identity section of $E_C/N_0=G_F$, lift recursively
through $E_C/N_{r+1}\to E_C/N_r$, whose kernels are finite.
The compatible sections give a continuous homomorphism
$G_F\to\varprojlim E_C/N_r=E_C$ splitting the projection.  This is an
actual section-existence proof, although it does not select a canonical
section.

For $F(T_1,\ldots,T_d)$ the arithmetic fundamental-group extension of
the constant curve is the pullback by restriction to $G_F$, so the
section pulls back continuously.  An $F(T_1,\ldots,T_d)$-point would
give a rational map $\mathbb P^d_F\dashrightarrow C$; restricting to
lines and using Riemann--Hurwitz shows that it is constant, and its
constant value is an $F$-point.  This contradicts $C(F)=\varnothing$.
\end{proof}

\paragraph{Statement D6 (Proved).}\label{res:D6}
The faithfulness assertion of manuscript Th\'eor\`eme 2 (scan 22,
manuscript p.~10) holds for finitely generated field extensions of
$F=K(\mu_\infty)$, including finite extensions of this $F$.
Explicitly, $F$-embeddings $f,g:M\hookrightarrow L$ inducing the same
outer continuous homomorphism $G_L\to G_M$ are equal.

\begin{proof}
As $\mu_n\subset F$, inner conjugation acts trivially on
$H^1(G_M,\mu_n)=\operatorname{Hom}_{\rm cont}(G_M,\mu_n)$.
The two outer maps therefore induce equal Kummer pullbacks for every
$n$, giving $f(a)/g(a)\in\bigcap_n L^{*n}$ for each $a\in M^*$.
Let $F_1$ be the relative algebraic closure of $F$ in $L$; it is finite
over $F$.  On a normal proper model of $L/F_1$, every prime-divisor
valuation of an infinitely divisible element is zero.  Such an element
and its inverse are global regular functions, hence constants in $F_1$.
Every root in $L$ of a constant is algebraic over $F_1$, so belongs to
$F_1$.  The arithmetic input in A(iv) now gives
$\bigcap_nL^{*n}=\mu_\infty$.

Choose an embedding of the countable characteristic-zero field $L$ into
$\mathbb C$.  Root-of-unity ratios give equal complex absolute values
for $f(a)$ and $g(a)$, and the same for $a+1$ and $a+i$.
The resulting equalities of squared absolute values give equal real
parts and equal imaginary parts of $f(a)$ and $g(a)$, respectively.
Here $f(i)=g(i)=i$ after naming the chosen image of $i$; the other image
$-i$ yields the same conclusion.  Thus $f(a)=g(a)$ for all $a$.
This argument also covers $a=-1,-i,0$ directly, since these belong to
$F$.  The comparison used throughout is the contravariant Kummer
pullback on $H^1$, so it respects the precise inner ambiguity in the
manuscript's groupoid formulation.
\end{proof}

\paragraph{Statement D7 (Refuted; exact replacement proved).}\label{res:D7}
The abelianized point-injectivity assertion of manuscript Th\'eor\`eme 3
(scan 32, manuscript p.~15) fails over $F$ and over its finitely generated
extensions.  Already for $\mathbb G_m/F$, the distinct points $1$ and
$\zeta_6$ have identical full section classes.  For an embedding of an
$F$-variety into a semiabelian variety, its induced Kummer classes agree
exactly when their difference belongs to the intersection of all
multiples of the rational-point group.  For hyperbolic curves the more
precise replacements are D1 and D2.

\begin{proof}
The geometric fundamental group of $\mathbb G_m$ is
$\widehat{\mathbb Z}(1)$, and the difference of the two point sections
is the Kummer class of their ratio.  Each root of unity has compatible
roots of all orders in $F$, so the class of $\zeta_6$ is zero.  This
persists under every extension of $F$.  For a semiabelian variety the
multiplication-by-$n$ exact sequence gives the stated divisibility
criterion at each finite level.  Compatible finite division-point sets
give the full Tate-module criterion, as in the proof of D1.
\end{proof}

\paragraph{Statement D8 (Proved).}\label{res:D8}
Manuscript Th\'eor\`eme 1 (scans 13--15, manuscript pp.~6--7), with its stated
centralizers of open subgroups, holds for every finitely generated
extension $M/F$.  It holds for its geometric absolute Galois group as
well, with the trivial group understood in transcendence degree zero.

\begin{proof}
The field $F$ is an abelian extension of a number field and is Hilbertian;
the exact Hilbertian inputs and the centralizer argument are proved in
Statement B12.  A finite extension preserves this property.  In positive
transcendence degree a rational function field is Hilbertian, and finite
extensions preserve it.  Apply B12 to $M$ and, for its geometric group,
to the positive-transcendence-degree function field over the algebraic
closure of the constant field.  This gives centralizer triviality for
every open subgroup, not only the centre.
\end{proof}

\paragraph{Statement D9 (Proved replacement; Conjecture A refuted).}\label{res:D9}
For the cyclotomic base $F$ and its finitely generated extensions, the
hyperbolic version of the Proposition of scan 54 and the centralizer
``pr\'ecisions'' of scan 56 hold with the same section and conjugacy
conventions as Parts A and C.  The excluded type $(0,2)$ remains a
counterexample to the first assertion of the Proposition.  Conjecture A
itself has the counterexample D5.

\begin{proof}
The finite-extension AVKF property used for packet separation holds by
D2.  The Tate-module vanishing for all finite extensions is the
arithmetic input proved in Part A/E; since the base contains all roots
of unity, a cyclotomically normalized inertia subgroup is centralized.
Apply the already proved inertia-recognition and unique-cusp argument
of Part A over this base and the splitting-class comparison of Part C.
Every comparison is a base extension of the stated finite \'etale
covering or the continuous inclusion of the corresponding open Galois
subgroup, so all its hypotheses persist.  The $(0,2)$ calculation is
literal: for $\mathbb G_m$ the decomposition extension at either cusp
is the full fundamental-group extension.  Its splitting packet
therefore contains every point section, including the pair in D7.
\end{proof}

\subsection{Poincar\'e coverings and the continuity clause}

\paragraph{Statement D10 (Proved).}\label{res:D10}
Let $k$ be an algebraically closed field of characteristic zero and let
$U,V$ be hyperbolic curves over $k$.  Their universal pro-\'etale
coverings are isomorphic as pro-curves over $k$ exactly when there is a
connected curve $W$ with finite \'etale surjections to both $U$ and $V$.
The completion of a universal covering is canonical and functorial for
semilinear isomorphisms.  More precisely, it is the inverse system of
smooth proper compactifications of its finite levels; the set used in
the fixed-point assertions is the set of compatible geometric points
\[
                       X'=\varprojlim_i X_i(k).
\]
Thus a generic scheme point is not being counted as a fixed point in
Conjectures B or B$'$.

\begin{proof}
The transition maps between finite levels are finite, hence affine.
Their scheme limit exists by Stacks, Lemma 32.2.2
\cite[Tag 01YV]{D_Stacks}.  A morphism from this limit to a curve of
finite presentation factors through a finite level.  The precise result
is Stacks, Proposition 32.6.1(3), Tag 01ZB: for an inverse system of
quasi-compact quasi-separated schemes with affine transitions, morphisms
to a locally finitely presented target are the filtered colimit of the
finite-level morphism sets.  The same statement identifies two such
morphisms after a sufficiently fine level when their limit morphisms
coincide.

Let $\alpha:\widetilde U\simeq\widetilde V$ be an isomorphism.  Factor
$p_V\alpha$ through a finite cover $U_i$ of $U$, obtaining $f:U_i\to V$.
Factor $p_{U_i}\alpha^{-1}$ through a finite cover $V_j$ of $V$, obtaining
$g:V_j\to U_i$.  After increasing $j$, the equality of limit morphisms
gives $fg=q_j$, where $q_j:V_j\to V$ is the finite \'etale structure
map.  In particular $f$ and $g$ are nonconstant.

Here is the finiteness step, which is necessary for the asserted
equivalence.  A nonconstant map of smooth curves is quasi-finite and
separated.  Zariski's main theorem, Stacks, Lemma 37.43.3, Tag 05K0,
factors $f$ as an open immersion $U_i\subset T$ followed by a finite map
$T\to V$; replace $T$ by the reduced closure of $U_i$ in $T$.
The composite $V_j\to T$ is proper over $V$: its graph is closed because
$T/V$ is separated, and projection from $V_j\times_V T$ is proper
because $V_j/V$ is finite.  Its image is closed and dense in $T$, so is
all of $T$, while it is contained in $U_i$.  Consequently $T=U_i$ and
$f$ is finite.  Since $g$ is surjective and $fg$ is \'etale, the
derivative of $f$ vanishes nowhere.  A finite map of smooth curves is
flat, as its local modules over the target discrete valuation rings
are torsion-free.  Hence $f$ is \'etale.  The curve $U_i$ is the required
common cover.

Conversely, the pointed finite \'etale covers of a connected finite
\'etale cover $W\to U$ form a cofinal subsystem of the covers of $U$:
for a given cover take the connected component, containing the chosen
geometric lift, of its fibre product with $W$.  Thus both $W\to U$ and
$W\to V$ identify their universal pro-coverings, with the usual deck
ambiguity when the lifts are forgotten.

A finite map between smooth curves extends uniquely between their
smooth proper models.  The exact algebraic statement is the equivalence
between transcendence-degree-one function fields and normal proper
curves with nonconstant morphisms, Stacks, Theorem 53.2.6, Tag 0BY1.
Apply it to the finite-level maps just constructed.  Equality after
refinement ensures compatibility, and the inverse isomorphism gives the
inverse map of completed systems.  This also proves independence of the
initial presentation, using the common cover.  Replacing $k$ by its
image under an automorphism gives the same construction for a semilinear
isomorphism.  On compatible geometric points these maps define exactly
the claimed functorial completion.
\end{proof}

\paragraph{Statement D11 (Proved).}\label{res:D11}
Let $K$ have characteristic zero and let $\mathcal U$ be a Poincar\'e
pro-curve over $\overline K$.  Give its semilinear automorphism group
the topology whose identity neighbourhoods are the simultaneous
stabilizers of finitely many rational functions.  Every semilinear
$G_K$-action for which each rational function has open stabilizer
preserves, after restriction to an open normal subgroup $G_{K_i}$, a
finite level $U_i$.  This level has a $K_i$-form, and the restricted
action is supplied by a continuous section of its arithmetic
fundamental-group extension.  One may start with any specified finite
level in this assertion.

The families of assertions B$'$ with this continuity clause and B with
$\Gamma^\circ$, each quantified over all finite extensions of $K$ and
all hyperbolic curves there, are equivalent.  At one fixed base $K$,
the precise implications are
\[
 \{\mathrm B_L:\ [L:K]<\infty\} \Longrightarrow\
 \mathrm B'_K\text{ (continuous)}\ \Longrightarrow\ \mathrm B_K.
\]
These are comparisons of the assertions, without assuming their
existence conclusions.  For the cyclotomic base in D5, even the
continuous form has a counterexample.

\begin{proof}
Let $\Omega$ be the union of the function fields of the finite levels,
and select one such field $L_0=\overline K(f_1,\ldots,f_r)$.
Choose equations for its smooth compactification and boundary, with
coefficients in a finite extension $K_0/K$.  Enlarge the finite list of
functions when necessary to realize these equations as the chosen
model.  Intersect the open stabilizers of the $f_j$ and the open
subgroup fixing $K_0$, and then take its normal core in $G_K$.
This gives an open normal $G_{K_i}$.  On $L_0$ its action fixes the
$f_j$ and acts on the coefficients in the prescribed way.  Thus it
is the standard semilinear action on that $K_i$-model, preserving its
boundary.  Equality of function-field actions implies equality of the
morphisms to this smooth separated curve.  The projection
$\mathcal U\to U_i$ is therefore equivariant.

The automorphisms of $\Omega$ fixing $K_i(U_i)$ form $E_{U_i}$: a
finite separable subextension is admitted precisely when its
normalization is \'etale over $U_i$, and adjoining the algebraic
constants gives the geometric universal cover.  This identifies the
actual group, including its projection to $G_{K_i}$, rather than merely
an outer action.  The given semilinear action therefore gives a
homomorphism $s_i:G_{K_i}\to E_{U_i}$ splitting that projection.
The Krull topology on this group is generated by stabilizers of finite
sets of elements of $\Omega$.  The hypothesis makes $s_i$ continuous.
In particular its image is compact.  Conversely a continuous section
has open stabilizers on every element of $\Omega$, since each element
lies in a finite arithmetic covering.  This proves both topological
comparisons directly.

Apply the fixed-point assertion B over $K_i$ to obtain its unique
$G_{K_i}^\circ$-fixed geometric pro-point.  The subgroup
$G_{K_i}^\circ=G_{K_i}\cap G_K^\circ$ is normal in $G_K$, so uniqueness
makes that point fixed by all of $G_K$.  Any $G_K^\circ$-fixed point
is also $G_{K_i}^\circ$-fixed, proving uniqueness in B$'_K$.
The reverse implication at a fixed base is the ordinary
section-to-action construction.  Applying these two implications at
every finite extension proves the displayed family equivalence.
Finally, the section of the proper curve in D5 has no fixed geometric
pro-point, because its image on the curve would be an $F$-rational
point.  Its action is continuous by the construction just proved.
\end{proof}

\paragraph{Statement D12 (Proved for proper curves).}\label{res:D12}
Let $U/\overline K$ be proper hyperbolic, choose an embedding
$\overline K\hookrightarrow\mathbb C$, and write
$U(\mathbb C)=\mathbb H/\Lambda$.  Its pro-curve automorphism group is
the Belyaev completion of $\operatorname{Comm}(\Lambda)$, with
$\widehat\Lambda$ a compact open subgroup.  In particular,
\[
 \operatorname{Aut}_{\overline K}(\widetilde U)
       =\widehat\Lambda\operatorname{Comm}(\Lambda),\qquad
 \operatorname{Aut}_{\overline K}(\widetilde U)/\widehat\Lambda
       \simeq\operatorname{Comm}(\Lambda)/\Lambda
\]
as groups in the first formula and as coset sets in the second.
For a non-arithmetic uniformizing group this completion is profinite;
for an arithmetic uniformizing group it is locally compact and
noncompact.

\begin{proof}
The opened source is D\v zambi\'c--Gonz\'alez-Diez
\cite[Theorem 7 and Proposition 3]{D_DzambicGonzalezDiez}.
Its hypothesis is a compact Riemann surface of genus at least two,
and its conclusion identifies the holomorphic automorphisms of the
inverse-limit solenoid with the completion of the commensurator in
the permutation action on cosets of its finite-index subgroups.
Proposition 3 states that $\widehat\Lambda$ is compact open and gives
the two displayed formulas.  The arithmeticity criterion used in its
Corollary 11 is
$[\operatorname{Comm}(\Lambda):\Lambda]=\infty$ precisely in the
arithmetic case; the source records this Margulis criterion explicitly
in its Introduction, p.~3.

Here is the comparison with the algebraic group used here.  Riemann
existence, SGA~1, XII, Theorem 5.1 \cite{D_SGA1}, sends each finite
topological cover to its finite \'etale algebraic cover, and finite
\'etale covers are unchanged by extension between algebraically closed
fields of characteristic zero (SGA~1, XIII, Corollary 2.9).
An algebraic pro-automorphism therefore acts holomorphically on the
inverse-limit solenoid.  Conversely, the source's formula expresses
every holomorphic automorphism as a deck element times a commensurator
element.  The latter gives an isomorphism between two finite covers
$\mathbb H/\Lambda_1$ and $\mathbb H/\Lambda_2$.  It is algebraic by
SGA~1, XII, Corollary 4.5, and
is defined over $\overline K$: the isomorphism scheme between two
proper genus-at-least-two curves is finite, so extending an algebraically
closed base cannot add its geometric points.  These two finite covers
and their isomorphism induce a pro-automorphism by the cofinality map
in D10.  This proves surjectivity of the comparison, while agreement
on finite-level function fields proves injectivity.  Stabilizers of
finite levels give the same compact open subgroups on the two sides.

When the commensurator has finite index over $\Lambda$, completing for
the finite-index-subgroup topology is its usual profinite completion.
When that index is infinite, the quotient by the compact open subgroup
$\widehat\Lambda$ is infinite discrete.  A compact group could not have
such a quotient.  This proves the two topological conclusions from
the stated criterion.
\end{proof}

\paragraph{Statement D13 (Open in the unrestricted cases; proved partial results).}
\label{res:D13}
The common-cover and completion assertions of scan 62, and the
continuity descent in Part D(e), are proved in D10--D11.  D12 gives the
automorphism description for proper curves.  An action over a number
field that normalizes no profinite deck group has not been constructed
or excluded here.  The affine commensurator-completion assertion, beyond
the finite-level comparisons of D10, has not been proved here either.
The smallest instance retained for this part is the Poincar\'e curve of
$\mathbb P^1\setminus\{0,1,\infty\}$ over $\overline{\mathbb Q}$,
with its full semilinear $G_{\mathbb Q}$-actions.

To state the missing assertion exactly, fix its usual continuous
tangential action $\rho_0$, put
$A=\operatorname{Aut}_{\overline{\mathbb Q}}(\mathcal U)$, and let
$\mathscr P$ be the profinite deck groups of finite algebraic curve
quotients of $\mathcal U$.  Every abstract semilinear action over
$G_{\mathbb Q}$ is uniquely a function $a:G_{\mathbb Q}\to A$ satisfying
\[
 a_{gh}=a_g\,\rho_0(g)a_h\rho_0(g)^{-1},
 \qquad \rho(g)=a_g\rho_0(g).
\]
The unresolved construction or exclusion asks whether such a function
can satisfy
\[
 \text{for every }P\in\mathscr P\text{ there exists }g\in G_{\mathbb Q}
 \text{ with }a_g\rho_0(g)P\rho_0(g)^{-1}a_g^{-1}\ne P.
\]
These are literal equations in the indicated automorphism group.
Their equivalence with the requested action follows by multiplying
$\rho(g)\rho(h)$ and by the definition of normalization; it does not
assume any continuity or any section conjecture.
For this Poincar\'e curve B$'$ asserts a unique
$G_{\mathbb Q}^\circ$-fixed compatible geometric point for every such
semilinear action.  D11 settles its continuous comparison with B;
the unrestricted assertion above is the outstanding part of D(e).
